\section{Exact normal spectral jets at integer inner orders}
\label{sp:sec:normal}

Write
\[
 \mathcal T(A,B,C;X,Y)
 =\sum_{m,n\geq1}\frac{X^mY^n}{m^A n^B(m+n)^C}.
\]
The orders are independent complex variables, and powers of positive
integers use real logarithms.  The question in the incoming
\emph{Resonance and Coordinate Transport} report asks for finite
reductions of the normal $C$-derivatives at $C=0$, beginning with
$A,B\in\{0,1\}$.  The negative-$C$ value identity by itself does not
answer that question.  The reduction below keeps $C$ variable from
the outset.

The underlying positive-integer partial-fraction identity is classical:
Zhao explicitly uses it to reduce colored Tornheim values to double
polylogarithms \cite{sp:Zhao}.  Our use of that identity retains the
complete complex $C$-dependence, specifies the continuation required
by the normal jets, and combines it with a rational-function theorem
covering arbitrary depth, integer weights, and coincident colors.
These are exact reductions to stated function classes.  No arithmetic
independence or global priority claim is inferred from them.

For the double polylogarithm use the convention
\begin{equation}\label{sp:doubleLi}
 \operatorname{Li}_{u,v}(z,w)
   =\sum_{k>n\geq1}\frac{z^k w^n}{k^u n^v}.
\end{equation}
Initially the colors can be taken with $|X|,|Y|<1$ and $X,Y\ne0$,
so that the reindexed series below converge locally normally for all
orders.  On $|X|=|Y|=1$ we assume $X,Y\ne1$ and use radial analytic
continuation: replace $(X,Y)$ by $(\rho X,\rho Y)$ and let
$\rho\uparrow1$.  In \eqref{sp:doubleLi} this means damping its outer
color by $\rho$.  This convention is unambiguous for the combinations
appearing here.  Indeed
\[
 \mathcal T(0,B,C;X,Y)
 =\frac1{\Gamma(C)}\int_0^\infty
 t^{C-1}\operatorname{Li}_0(Xe^{-t})
                    \operatorname{Li}_B(Ye^{-t})\,dt
\]
for $\Re C>0$.  The product is analytic at $t=0$ and exponentially
decreasing at infinity.  Taylor subtraction at zero, followed by the
zeros of $1/\Gamma(C)$, gives an entire continuation in $(B,C)$.
The same argument applies with arbitrary fixed integer first orders,
and identifies the radial continuation locally uniformly in the
orders.  Thus all the order derivatives used below exist.

\begin{theorem}[Integer inner orders and all normal jets]
\label{sp:integer}
For positive integers $A,B$, one has the identity of entire functions
of $C$ at nontrivial unit colors
\begin{align}\label{sp:positive}
 \mathcal T(A,B,C;X,Y)
 ={}&\sum_{j=1}^{A}\binom{A+B-j-1}{B-1}
       \operatorname{Li}_{C+A+B-j,j}(Y,X/Y)\notag\\
 &+\sum_{j=1}^{B}\binom{A+B-j-1}{A-1}
       \operatorname{Li}_{C+A+B-j,j}(X,Y/X).
\end{align}
For $M\geq0$ an integer and arbitrary $B\in\mathbb C$,
\begin{equation}\label{sp:nonpositive}
 \mathcal T(-M,B,C;X,Y)
 =\sum_{j=0}^{M}(-1)^j\binom Mj
       \operatorname{Li}_{C-M+j,B-j}(X,Y/X).
\end{equation}
Every fixed $C$-derivative can be applied to these formulas.  In
\eqref{sp:nonpositive}, every fixed $B$-derivative is also valid.
Together with exchange of the two inner variables, these formulas
give a finite reduction for every pair of integer inner orders.
\end{theorem}
\begin{proof}
For positive $m,n$,
\[
 \frac1{m^A n^B}
 =\sum_{j=1}^{A}\binom{A+B-j-1}{B-1}
     \frac1{m^j(m+n)^{A+B-j}}
 +\sum_{j=1}^{B}\binom{A+B-j-1}{A-1}
     \frac1{n^j(m+n)^{A+B-j}}.
\]
It follows by repeated use of
$1/(mn)=(1/m+1/n)/(m+n)$, or by induction on $A+B$.
Multiply by the colors and $(m+n)^{-C}$ and put $k=m+n$.
The first sum has color $Y^k(X/Y)^m$, and the second has
$X^k(Y/X)^n$.  This proves \eqref{sp:positive} in an absolutely
convergent region.  For \eqref{sp:nonpositive}, expand
$(k-n)^M=\sum_j(-1)^j\binom Mj k^{M-j}n^j$ after the same reindexing.
Local normal convergence inside the bidisc and then the specified
analytic continuation prove both identities and the derivative claims.
\end{proof}

Here is the promised complete $A,B\in\{0,1\}$ table.  A superscript
$[r]$ on a double polylogarithm in this table differentiates only its
first order.  For $X\ne Y$,
\begin{align}
 \partial_C^r\mathcal T(0,0,C;X,Y)
 &=\frac{Y\operatorname{Li}_{C}^{[r]}(X)
          -X\operatorname{Li}_{C}^{[r]}(Y)}{X-Y},
 \label{sp:00}\\
 \partial_C^r\mathcal T(0,1,C;X,Y)
 &=\operatorname{Li}_{C,1}^{[r]}(X,Y/X),\label{sp:01}\\
 \partial_C^r\mathcal T(1,0,C;X,Y)
 &=\operatorname{Li}_{C,1}^{[r]}(Y,X/Y),\label{sp:10}\\
 \partial_C^r\mathcal T(1,1,C;X,Y)
 &=\operatorname{Li}_{C+1,1}^{[r]}(X,Y/X)
   +\operatorname{Li}_{C+1,1}^{[r]}(Y,X/Y).
 \label{sp:11}
\end{align}
The confluent version of \eqref{sp:00} is
\begin{equation}\label{sp:00confluent}
 \partial_C^r\mathcal T(0,0,C;z,z)
 =\operatorname{Li}_{C-1}^{[r]}(z)
    -\operatorname{Li}_{C}^{[r]}(z).
\end{equation}
All $C=-N$ and $C=0$ substitutions are now legitimate.  Formula
\eqref{sp:11} is a finite depth-two answer; it does not claim a
reduction of the remaining depth-two order derivatives to ordinary
polylogarithms.

\begin{remark}[The permitted derivative directions]
\label{sp:directionwarning}
In \eqref{sp:positive}, the first two orders are fixed integers.  Its
finite expression does not authorize differentiating with respect to
$A$ or $B$.  In particular, it does not reduce the full mixed jet
$\partial_A\partial_B\partial_C\mathcal T$ by formal differentiation
of the finite sums.  In \eqref{sp:nonpositive}, $B$ really is free,
so $B$-derivatives are allowed, while $A$ remains the fixed integer
$-M$.  This is the same distinction that prevents differentiating
a negative-$C$ value formula transversely without further information.
\end{remark}

\subsection{Rational generating functions give depth-one closure}

Let $d\geq1$, let $p_j$ be positive integers, let $a_j$ be nonnegative
integers, and let $X_j$ be nonzero complex numbers with $|X_j|\leq1$.
Define, initially for $\Re C>a_1+\cdots+a_d+d$,
\begin{equation}\label{sp:weighteddef}
 Z(C)=\sum_{m_1,\ldots,m_d\geq1}
  \frac{\prod_{j=1}^d m_j^{a_j}X_j^{m_j}}
       {(p_1m_1+\cdots+p_dm_d)^C}.
\end{equation}
Consider its coefficient generating function
\begin{equation}\label{sp:rationalR}
 R(t)=\prod_{j=1}^d\operatorname{Li}_{-a_j}(X_jt^{p_j})
      =\sum_{n\geq1}r_nt^n.
\end{equation}
The elementary identity
$\operatorname{Li}_{-a}(z)=(z\,d/dz)^a\{z/(1-z)\}$ proves that
$R$ is rational and bounded at infinity.  Write its partial fractions
in the form
\begin{equation}\label{sp:partialR}
 R(t)=c_\infty+
       \sum_{\alpha}\sum_{h=1}^{M_\alpha}
           \frac{c_{\alpha,h}}{(1-\alpha t)^h},
\end{equation}
where each $\alpha$ is a root of some equation
$\alpha^{p_j}=X_j$.  A permissible multiplicity bound is
\[
 M_\alpha=\sum_{j:\,\alpha^{p_j}=X_j}(a_j+1).
\]
This is an upper bound: zeros of other factors can cancel some poles.
In that event the corresponding highest coefficients are zero.
There is no positive-degree polynomial part because $R$ is bounded
at infinity.  An exact coefficient formula, also valid at repeated
roots, is
\begin{equation}\label{sp:cformula}
 c_{\alpha,h}
   =[u^{M_\alpha-h}]\,
       u^{M_\alpha}R((1-u)/\alpha).
\end{equation}
The function whose coefficient is extracted is holomorphic at $u=0$.

\begin{theorem}[Arbitrary depth, weights, multiplicities, and normal jets]
\label{sp:rationaltheorem}
Set
\begin{align}
 e_{h,r}
  &=[u^r]\prod_{\nu=1}^{h-1}(1+u/\nu),\label{sp:eh}\\
 d_{\alpha,r}
  &=\sum_{h=r+1}^{M_\alpha}c_{\alpha,h}e_{h,r}.
 \label{sp:dcoeff}
\end{align}
Empty products equal one.  Then
\begin{equation}\label{sp:depthone}
 Z(C)=\sum_\alpha\sum_{r=0}^{M_\alpha-1}
               d_{\alpha,r}\operatorname{Li}_{C-r}(\alpha).
\end{equation}
This is the meromorphic continuation of \eqref{sp:weighteddef}, and
every $C$-derivative is obtained by replacing each polylogarithm by
the corresponding order derivative.

If no $X_j$ equals one, the continuation is entire.  In general its
only possible poles are simple poles at $C=r+1$ from the $\alpha=1$
sector, with residue $d_{1,r}$.  For $k\geq0$, the constant Laurent
coefficient of $Z^{(k)}$ at an integer $m\geq1$ is
\begin{equation}\label{sp:stieltjesFP}
 \operatorname{CT}_{C=m} Z^{(k)}(C)
 =(-1)^k d_{1,m-1}\gamma_k
   +\sum_{(\alpha,r)\ne(1,m-1)}
      d_{\alpha,r}\operatorname{Li}_{m-r}^{[k]}(\alpha).
\end{equation}
Missing coefficients are zero; for $\alpha=1$ the nonsingular terms
mean ordinary Riemann-zeta derivatives.
\end{theorem}
\begin{proof}
The original multiple sum converges absolutely in the stated
half-plane.  Grouping it according to $n=\sum p_jm_j$ gives
$Z(C)=\sum_{n\geq1}r_n n^{-C}$.  The binomial series gives
\[
 [t^n](1-\alpha t)^{-h}
 =\alpha^n\binom{n+h-1}{h-1}
 =\alpha^n\prod_{\nu=1}^{h-1}(1+n/\nu)
 =\alpha^n\sum_{r=0}^{h-1}e_{h,r}n^r.
\]
For $n\geq1$ the constant $c_\infty$ contributes nothing.  Substitution
proves \eqref{sp:depthone} where the series converges absolutely and hence
by continuation.  If $\alpha\ne1$ and $|\alpha|\leq1$, the fixed-color
polylogarithm is entire in its order; for $\alpha=1$ it is $\zeta$.
The pole and derivative statements follow.  Finally
\[
 \zeta(1+v)=v^{-1}+
       \sum_{j\geq0}\frac{(-1)^j\gamma_j}{j!}v^j
\]
shows that the constant coefficient after $k$ differentiations is
$(-1)^k\gamma_k$.  The other terms are holomorphic at the specified
integer, proving \eqref{sp:stieltjesFP}.
\end{proof}

The harmonic-number coefficients in this theorem can be computed by
the finite-degree expansion
\begin{equation}\label{sp:harmonicexp}
 \prod_{\nu=1}^{h-1}(1+u/\nu)
 =\exp\!\left(
    \sum_{j\geq1}\frac{(-1)^{j+1}}j H_{h-1}^{(j)}u^j\right).
\end{equation}
For example $e_{h,0}=1$, $e_{h,1}=H_{h-1}$, and
$e_{h,2}=(H_{h-1}^2-H_{h-1}^{(2)})/2$.
All multiplicity corrections are consequently finite polynomials in
generalized harmonic numbers.  Coincidence in this statement concerns
the poles of the rational coefficient generating function.  It does
not identify two independently specified spatial finite parts.

\subsection{Roots of unity, Gamma values, and an explicit weighted example}

For a nontrivial $q$th root of unity $\alpha$, residue-class summation
gives
\begin{equation}\label{sp:DFT}
 \operatorname{Li}_s(\alpha)
   =q^{-s}\sum_{j=1}^q\alpha^j\zeta(s,j/q).
\end{equation}
It follows first where the series converge absolutely and then by
continuation.  In particular, for $s\ne1$,
\begin{equation}\label{sp:DFTjets}
 \operatorname{Li}_s^{[k]}(\alpha)
 =q^{-s}\sum_{j=1}^q\alpha^j
       \sum_{\ell=0}^k\binom{k}{\ell}(-\log q)^{k-\ell}
                 \zeta^{[\ell]}(s,j/q).
\end{equation}
At $s=1$, cancellation of the common pole gives instead
\begin{equation}\label{sp:DFTstieltjes}
 \operatorname{Li}_1^{[k]}(\alpha)
 =\frac1q\sum_{j=1}^q\alpha^j
       \sum_{\ell=0}^k\binom{k}{\ell}
        (-\log q)^{k-\ell}(-1)^\ell\gamma_\ell(j/q).
\end{equation}
The cancellation is performed before evaluation.  Applying Lerch's
formula $\zeta'(0,a)=\log\Gamma(a)-\tfrac12\log(2\pi)$
\cite[Eq.~25.11.18]{DLMFHzeta} gives
\begin{equation}\label{sp:Gammajet}
 \operatorname{Li}_0^{[1]}(\alpha)
   =\sum_{j=1}^q\alpha^j\log\Gamma(j/q)
       -\log q\,\frac{\alpha}{1-\alpha}.
\end{equation}
Thus simple poles in \eqref{sp:partialR} yield a finite log-Gamma
evaluation for $Z'(0)$ at roots of unity.  Repeated poles additionally
introduce the explicitly listed Hurwitz derivatives at negative
integers from \eqref{sp:DFTjets}; they should not be silently discarded.

For a concrete unequal-weight example, put
\[
 Z_{2,3}(C)=\sum_{m,n\geq1}\frac{(-1)^{m+n}}{(2m+3n)^C}
 \qquad(\Re C>2).
\]
Its rational generating function has the exact certificate
\begin{equation}\label{sp:weightedcertificate}
 \frac{t^5}{(1+t^2)(1+t^3)}
 =\frac{t^5-t^7-t^8+t^9+t^{10}-t^{12}}{1-t^{12}}.
\end{equation}
Consequently, for every $C$ by entire continuation,
\begin{align}\label{sp:weightedHurwitz}
 Z_{2,3}(C)=12^{-C}\bigl\{
  &\zeta(C,5/12)-\zeta(C,7/12)-\zeta(C,2/3)\notag\\
  &+\zeta(C,3/4)+\zeta(C,5/6)-\zeta(C,1)\bigr\}.
\end{align}
The six residues at $C=1$ cancel.  In particular,
\begin{equation}\label{sp:weightedGamma}
 Z_{2,3}'(0)
 =\log\frac{\Gamma(5/12)\Gamma(3/4)\Gamma(5/6)}
                 {\Gamma(7/12)\Gamma(2/3)}-\frac14\log12.
\end{equation}
An entirely convergent integral version, with no formal double sum
at $C=0$, is
\begin{align}\label{sp:weightedIntegral}
 &\int_0^1\left\{\frac1{(e^{2t}+1)(e^{3t}+1)}-\frac14\right\}
                         \frac{dt}{t}
 +\int_1^\infty\frac{dt}{t(e^{2t}+1)(e^{3t}+1)}
 +\frac\gamma4\notag\\
 &\hspace{24mm}=
 \log\frac{\Gamma(5/12)\Gamma(3/4)\Gamma(5/6)}
                 {\Gamma(7/12)\Gamma(2/3)}-\frac14\log12.
\end{align}
To verify the integral, insert the Mellin representation for
$\Re C>0$, subtract its constant $1/4$ at zero, and use
$1/\Gamma(C)=C+\gamma C^2+O(C^3)$.

For a fully confluent example, $d$ identical colors with unit weights
and zero inner orders give
\begin{equation}\label{sp:compositions}
 \sum_{m_1,\ldots,m_d\geq1}
       \frac{z^{m_1+\cdots+m_d}}{(m_1+\cdots+m_d)^C}
 =\frac1{(d-1)!}\sum_{r=0}^{d-1}s(d,r+1)
                         \operatorname{Li}_{C-r}(z),
\end{equation}
where $s(d,j)$ are the signed Stirling numbers of the first kind.
Indeed the coefficient at total index $n$ counts compositions and
equals $\binom{n-1}{d-1}$, including its zero values for $1\leq n<d$.
This proves the identity and all its normal derivatives.  For $d=3$,
$z=-1$, and differentiation at $C=0$, it yields
\begin{equation}\label{sp:triplealternating}
 \left.\frac{d}{dC}
 \sum_{m,n,k\geq1}\frac{(-1)^{m+n+k}}{(m+n+k)^C}
 \right|_{C=0}^{\mathrm{cont}}
 =-\frac{7\zeta(3)}{8\pi^2}
   -\frac92\zeta'(-1)-\frac12\log\pi.
\end{equation}
Here ``cont'' explicitly means continuation from $\Re C>3$.
The reduction uses
$\operatorname{Li}_s(-1)=(2^{1-s}-1)\zeta(s)$,
$\zeta'(0)=-\tfrac12\log(2\pi)$, and
$\zeta'(-2)=-\zeta(3)/(4\pi^2)$.

\subsection{An exact primitive hierarchy}

Keep the coefficients $r_n$ of \eqref{sp:rationalR} and put
$Z(C;z)=\sum_{n\geq1}r_n z^n/n^C$ for $|z|<1$.
The proof of Theorem~\ref{sp:rationaltheorem} gives
\[
 Z(C;z)=\sum_{\alpha,r}d_{\alpha,r}
                    \operatorname{Li}_{C-r}(\alpha z).
\]
For every nonnegative integer $k$,
\begin{equation}\label{sp:primitive}
 \int_0^z Z^{(k)}(C;t)\frac{dt}{t}=Z^{(k)}(C+1;z),
 \qquad
 z\frac{d}{dz}Z^{(k)}(C;z)=Z^{(k)}(C-1;z).
\end{equation}
The integrand is holomorphic at zero after division by $t$, and the
primitive is fixed by its zero value there.  These identities follow
termwise in the unit disc and continue along any common path avoiding
the polylogarithmic singularities.  This gives exact antiderivatives
of every normal jet without changing its spectral direction.

\subsection{Independent algebraic and numerical checks}

The accompanying script \texttt{verify\_spectral.py} performs 402
exact finite checks: 25 positive-order partial fractions, 12
composition polynomials, 105 harmonic--Stirling coefficients, seven
reconstructions of rational generating functions, 252 independently
convolved coefficients, and the cyclotomic certificate
\eqref{sp:weightedcertificate}.  The rational cases include unequal
weights, repeated roots, multiplicity six, and cancellation of an
otherwise apparent pole.  Nine independently evaluated convergent
Mellin integrals at 60 decimal digits check the first normal jets for
distinct colors $i,-i$, repeated color $-1$ in depths two and three,
the unequal weights $(2,3)$, and the root-Gamma formula for roots of
orders $3,4,5,6,12$.  The largest observed discrepancy is
$4.7\times10^{-61}$.  The exact algebraic checks and analytic proofs
are separate from these floating-point diagnostics; the latter are
not interval-certified error bounds.  Inputs, formulas, and full
outputs are supplied in \texttt{spectral\_results.json}.
